\section{Intervalos disjuntos (ODT)}
\whenuse{El arreglo est\'a formado por pocos intervalos constantes y predominan
asignaciones de rango; apropiado con fragmentaci\'on controlada o datos aleatorios.}

Usar para mantener una partici\'on ordenada en intervalos constantes y aplicar
asignaciones sobre rangos. Funciona muy bien con datos aleatorios o cuando cada
operaci\'on crea pocos cortes; no garantiza buen peor caso.

\complexity{\texttt{split} $O(\log I)$; una operaci\'on cuesta
$O((I_{borrados}+1)\log I)$; memoria $O(I)$. Peor caso total cuadr\'atico.}

\begin{lstlisting}[style=cpp]
struct DisjointIntervals {
    struct Interval { int right; mutable ll value; };
    map<int, Interval> intervals; // left -> {right,value}, rangos [left,right)

    auto split(int position) {
        auto iterator = prev(intervals.upper_bound(position));
        if (iterator->first == position) return iterator;
        int right = iterator->second.right;
        ll value = iterator->second.value;
        if (position >= right) return intervals.end();
        iterator->second.right = position;
        return intervals.emplace(position, Interval{right, value}).first;
    }

    void assign(int left, int right, ll value) {
        auto after = split(right);
        auto first = split(left);
        intervals.erase(first, after);
        intervals[left] = {right, value};
    }
};
odt.assign(100, 500, 1);
odt.assign(300, 800, 2);
odt.assign(700, 900, 3);
\end{lstlisting}

\warning{Inicializar la partici\'on cubriendo todo el dominio. No usar ODT si un
adversario puede fragmentar los intervalos en cada operaci\'on.}
